\hspace{3mm}

\centerline{\bf \huge 11 класс}
\centerline{\normalsize{\it Работа рассчитана на \bf{210} минут}}

\hspace{3mm}

\problem{11.1}Аюка выписал на доску многочлен $19x^2 - 2023x + 2005$, а Мингиян возвёл этот многочлен в $2023$ степень, раскрыл скобки, привёл подобные и посчитал сумму коэффициентов этого многочлена. Какое число у него получилось?

\hspace{3mm}

\problem{11.2}Алтан и Бамба записали по одному целому числу на доску и рядом написали сумму этих чисел. Затем Алтан умножил своё число на $X$, а число Бамбы на $Y$, сложил полученные произведения и оказалось, что новая сумма делится на сумму чисел, записанную ребятами в самом начале на доске. Бамба решил поступить по-своему и умножил своё число на $X$, а число Алтана на $Y$ и сложил полученные произведения. Докажите, что новая сумма Бамбы также делится на изначальную сумму, записанную ребятами на доске, если известно, что $X$, $Y$ - целые числа.

\hspace{3mm}

\problem{11.3}Джангар каждый будний день недели побеждает некоторое количество противников, а в выходные дни он отдыхает. Известно, что в каждый последующий день недели он побеждал не больше противников, чем в предыдущий. И за всю третью неделю декабря, начиная с понедельника, он победил 6000 врагов. Какое наименьшее суммарное количество врагов он мог победить за первый, третий и пятый дни в третью неделю декабря? 

\hspace{3mm}

\problem{11.4}Известно что $a, b, c \geq 0$
Докажите, что $(8a + b + c)(a + \dfrac{1}{8}b + c)(a + b + \dfrac{1}{27} c ) \geq 9abc$

\hspace{3mm}

\problem{11.5}Дан остроугольный треугольник ABC. Окружности с центрами $A$ и $C$ проходят через точку
$B$, вторично пересекаются в точке $F$ и пересекают описанную около треугольника ABC окружность $\omega$
в точках $D$ и $E$. Отрезок $BF$ пересекает окружность $\omega$ в точке $O$. Докажите, что $O$ — центр описанной
окружности треугольника $DEF$.
